\section{智能的衡量：生存、探索、自动化与 token 经济}

上一章给出路线图，本章回答张小珺追问的“智能本质是什么”。广密没有给出严格定义，而是用三个关键词描述人类智能进化：生存、探索、自动化。放到模型上，智能进步可以看作模型在环境中更好地适应、发现新策略、自动完成任务，并用更低或更有效的 token 成本完成更复杂工作。

\begin{figure}[H]
\centering
\includegraphics[width=0.92\textwidth]{intelligence-measurement.png}
\caption{智能进步如何衡量：从生存、探索、自动化，到 token 消耗和任务闭环。自制概念图，依据 00:43:00--00:48:03 对谈内容整理。}
\end{figure}

\begin{knowledgebox}{读图：智能不是一个 benchmark 数字}
图中有生存、探索、自动化、Token、任务和反馈。读这张图时要注意：benchmark 是测量切片，真实智能更像任务闭环。模型能否独立探索、选择工具、消耗合理 token、完成可验证任务，是更接近 Agent 时代的衡量方式。
\end{knowledgebox}

\subsection{Token 消耗为什么重要}

本节解释 token 经济。访谈中提到，一个普通 chatbot 对话可能消耗几千 token，一个搜索或研究任务可能消耗更多，而 Manus 平均任务可能达到几十万 token。数字不必被当作精确统计，关键是趋势：Agent 任务比聊天任务消耗更多推理、上下文、工具调用和中间状态。产品定价如果仍然照搬 SaaS 月费，可能会和真实计算成本错配。

\begin{importantbox}{实践经验：AI 产品定价必须理解 token 成本}
传统 SaaS 的边际成本较低，而 Agent 可能每次任务都消耗大量 token、工具调用和外部计算。如果产品按固定月费卖，但用户大量运行高成本任务，商业模式会被 token 经济反噬。
\end{importantbox}

\subsection{从 long context 到 long-term memory}

本节为后面的 Agent 做铺垫。广密提到，AGI 接下来的 milestone 可能是 long-term memory，它会取代单纯 long context。Long context 是把更多材料塞进当前窗口；long-term memory 是模型能持续积累、选择、更新和调用历史经验。前者解决“这一轮看得更多”，后者解决“长期越用越懂、越做越会”。

\begin{knowledgebox}{术语消化：Long Context 与 Long-term Memory}
\noindent\begin{tabularx}{\textwidth}{p{0.22\textwidth}p{0.34\textwidth}X}
\toprule
概念 & 含义 & Agent 中的作用 \\
\midrule
Long Context & 单次推理窗口可容纳更长材料 & 让模型一次性读更多文件、网页和历史。 \\
Long-term Memory & 跨会话、跨任务保存和更新经验 & 让模型形成用户、项目和环境的持续状态。 \\
Task State & 当前任务的目标、步骤、错误和产物 & 决定 Agent 能否长程执行。 \\
Experience Reuse & 把过去成功/失败迁移到新任务 & 是 Online Learning 的前置能力之一。 \\
\bottomrule
\end{tabularx}
\end{knowledgebox}

\subsection{本章小结}

智能衡量不能只看单题分数。Agent 时代更重要的是任务闭环、环境反馈、token 成本、长期记忆和经验复用。这些变量共同决定模型是否真的从“会答”走向“会做”。

